% \section{Metatheory}
% \label{sec:metatheory}

% This section introduces the main results of our language: preservation and
% progress possibly ending in a deadlock.

\begin{figure}[t!]
  \declrel{Liveness/context agreement}
    {$\operator{LiveCtx}(L;\Gamma)$}
  \begin{equation*}
    \operator{LiveCtx}(L;\Gamma)
    \quad\Longleftrightarrow\quad
    \operator{ends}(L)=\operator{dom}(\Gamma),
    \qquad
    \operator{ends}(L)=\{x\mid\exists y.\ \{x,y\}\in L\}.
  \end{equation*}

  \declrel{Splitting a configuration context}
    {$\Gamma=\Gamma_1+\Gamma_2$}
  \begin{mathpar}
    \axiom{S-Empty}
    {\Empty=\Empty+\Empty}

    \infrule{S-LinL}
    {\Gamma=\Gamma_1+\Gamma_2}
    {\Gamma,\ebind xS=(\Gamma_1,\ebind xS)+\Gamma_2}

    \infrule{S-LinR}
    {\Gamma=\Gamma_1+\Gamma_2}
    {\Gamma,\ebind xS=\Gamma_1+(\Gamma_2,\ebind xS)}
  \end{mathpar}

  \declrel{Typing a thread multiset}{$\Gamma\vdash_{\mathsf{th}}E$}
  \begin{mathpar}
    \axiom{TM-Empty}
    {\Empty\vdash_{\mathsf{th}}\emptyset}

    \infrule{TM-Thread}
    {\Gamma=\Gamma_e+\Gamma_r \\
     \expAgainst[\Empty]{\Gamma_e}{e}{\TUnit}{\Empty} \\
     \Gamma_r\vdash_{\mathsf{th}}E}
    {\Gamma\vdash_{\mathsf{th}}\{e\}\uplus E}
  \end{mathpar}

  \declrel{Typing a process configuration}{$\isProc K$}
  \begin{mathpar}
    \infrule{T-Conf}
    {\operator{LiveCtx}(L;\Gamma) \and
     \Gamma\vdash_{\mathsf{th}}E}
    {\isProc[\Gamma]{(E,L)}}
  \end{mathpar}
  \caption{Typing rules for process configurations.  Context addition is
    disjoint union, and $\operator{dom}(\Gamma)$ is the set of variables bound
    in $\Gamma$.}
  \label{fig:proc-formation}
\end{figure}

%%% Local Variables:
%%% mode: latex
%%% TeX-master: "main"
%%% End:


We complete the section with the main results of our language.
%
Towards this end we need to talk about context and configuration typing.
%
Type formation (\cref{fig:type-formation}) is lifted pointwise to contexts in 
judgement $\isCtx{\Gamma}{\kind}$.
%
% , defined by the rules below. Notice that the
% last premise in rule Ctx-Empty checks the invariant whereby contexts contain
% only normalized terms.
% %
% \begin{mathpar}
%      \infrule{Ctx-Empty}
%      {}
%      {\isCtx{\Empty}{\kind}}
%  
%      \infrule{Ctx-Ext}
%      {\isCtx{\Gamma}{\kind}\\ \typeAgainst{T}{\kind} \\ \typec T = \Normal[+]{T}}
%      {\isCtx{\Gamma, \ebind{x}{T}}{\kind}}
% \end{mathpar}
%
The rules for typing process configurations are in
\cref{fig:proc-formation}.  A configuration context records named channel
endpoints and their linear session types.  Thus
$\operator{LiveCtx}(L;\Gamma)$ states that the domain of $\Gamma$ is exactly
the set of endpoints occurring in the pairs of $L$.  Endpoints of closed
channels do not occur in the context.  The judgment
$\operator{PairedCtx}(L;\Gamma)$ requires the two endpoint types of each live
channel to be normalized duals.  Because channel pairs are unordered and
duality is involutive, this condition is independent of the order in which
their endpoints are named.
%
The thread-multiset judgment selects an expression $e$ from a decomposition
$\mset{e}\uplus E$ and partitions the context as
$\Gamma=\Gamma_e+\Gamma_r$.  Context addition is disjoint, so each binding is
assigned either to $e$ or to the remainder $E$.  The selected expression must
check against $\TUnit$ and consume its entire context, leaving the empty
context.  Correspondingly, the empty thread multiset is typable only under the
empty context.  Consequently every live endpoint is owned and consumed by
exactly one thread.

Expression preservation still uses the labelled context transition from
\cref{sec:proofs} (\cref{fig:lts-ctx}).  At the configuration level the updated
context is instead constructed together with the target typing derivation.
The paired-context invariant supplies the compatibility of the endpoints
selected by a synchronization.  For value communication, preservation
transfers exactly the resources used to type the payload from the sender's
context share to the receiver's and reconstructs the target split.  Selection
and closing similarly use endpoint duality to advance or remove both endpoint
types.  The mechanization contains the corresponding resource-extraction and
context-reconstruction lemmas.

\begin{theorem}[Preservation]\
  \label{thm:preservation}
  \begin{enumerate}
  \item\label{it:pres-synth} If $\trans{e_1}{\lambda}{e_2}$ and
    $\expSynth{\Gamma_1,\Gamma_2,\Gamma_3}{e_1}{T}{\Gamma_3}$ and
    $\transCtx{\Gamma_2}{\lambda}{\Gamma_2'}$ and 
    $\isCtx{\Gamma_1,\Gamma_2,\Gamma_3}{\kindt}$, then
    $\expSynth{\Gamma_1,\Gamma_2',\Gamma_3}{e_2}{T'}{\Gamma_3}$ and
    $\typec{T'} \subt \typec{T}$.
  \item\label{it:pres-against} If $\trans{e_1}{\lambda}{e_2}$ and
    $\expAgainst{\Gamma_1,\Gamma_2,\Gamma_3}{e_1}{T}{\Gamma_3}$ and
    $\transCtx{\Gamma_2}{\lambda}{\Gamma_2'}$ and 
    $\isCtx{\Gamma_1,\Gamma_2,\Gamma_3}{\kindt}$, then
    $\expAgainst{\Gamma_1,\Gamma_2',\Gamma_3}{e_2}{T}{\Gamma_3}$.
  \item\label{it:pres-proc} If $\isProc[\Gamma]K$ and
    $\trans K\pi{K'}$, then there exists a context $\Gamma'$ such that
    $\isProc[\Gamma']{K'}$.
  \end{enumerate}
\end{theorem}

%\paragraph{Progress}

For progress, expression actions are divided according to whether they need a
second thread.  An expression is \emph{runnable} if it can take a $\beta$,
fork, or allocation step.  It is \emph{communication-blocked} if it has a
value-receive, value-send, selector-receive, selector-send, or close transition;
such an expression is not stuck locally, but its transition needs a matching
thread.  We write $\operator{RunnableAt}(K)$ when some thread is
runnable and $\operator{Sync}(K)$ when two distinct thread occurrences satisfy
the premises of Act-Msg, Act-Bra, or Act-Wait.

A configuration is \emph{terminal} if every thread is a value.  A
\emph{global deadlock} is positive evidence consisting of all four properties
below:
%
\begin{enumerate}
\item every thread is either a value or communication-blocked;
\item at least one thread is communication-blocked;
\item $\operator{RunnableAt}(K)$ does not hold; and
\item $\operator{Sync}(K)$ does not hold.
\end{enumerate}
%
The last two clauses exclude every source of a configuration transition, so a
global deadlock cannot step.

Local progress holds for expressions in session-only run-time contexts: an
expression that checks against $\TUnit$ is a value, is runnable, or exposes a
communication action.  Moreover, both $\operator{RunnableAt}(K)$ and
$\operator{Sync}(K)$ are decidable by finite search.  Lifting the local result
over the typed thread multiset and applying these decisions yields the
following trichotomy.

\begin{theorem}[Configuration progress]\
  \label{thm:progress}
  If $\isProc[\Gamma]K$, then either $\termc K$ is terminal, it is a global
  deadlock, or there exist $\pi$ and $\termc{K'}$ such that
  $\trans K\pi{K'}$.
\end{theorem}
%

Write $\termc K\Longrightarrow^*\termc{K'}$ for the reflexive-transitive
closure of configuration reduction, with labels omitted.  Iterating
preservation and applying configuration progress at the endpoint gives finite
type safety.

\begin{theorem}[Finite type safety]\
  \label{thm:finite-safety}
  If $\isProc[\Gamma]K$ and
  $\termc K\Longrightarrow^*\termc{K'}$, then there exists a context
  $\Gamma'$ such that $\isProc[\Gamma']{K'}$.  Moreover, $\termc{K'}$ is
  terminal, is a global deadlock, or can take a configuration step.
\end{theorem}

If a closed expression $e$ checks against $\TUnit$ from and to the empty
context, then $(\mset{e},\emptyset)$ is well typed.  Thus every finite reduction
from this singleton configuration ends in a terminal state, a global
deadlock, or a state that can step, including after channel allocations.


%%% Local Variables:
%%% mode: latex
%%% TeX-master: "main"
%%% End:
